%-------------------------------------------------------
%	CLASS, PACKAGES AND OTHER DOCUMENT CONFIGURATIONS
%-------------------------------------------------------
\documentclass{article}[
	a4paper, % Paper size
	12pt,    % Font size
	%unnumberedsections, % Uncomment for no section numbering
]
%-------------------------------------------------------
%	CLASS CONFIGURATION
%-------------------------------------------------------
\usepackage{etoolbox} % Required for conditional logic and easily changing commands

\newtoggle{unnumberedsections} 			% Create toggle for a class option
\settoggle{unnumberedsections}{false} 	% Default value for the class option
\DeclareOption{unnumberedsections}{\settoggle{unnumberedsections}{true}} % Set the class option toggle if the class option was used in the template

\DeclareOption*{\PassOptionsToClass{\CurrentOption}{article}} 			 % Pass through any extra options specified to the base class
\ProcessOptions\relax 					% Process class options

%-------------------------------------------------------
%	REQUIRED PACKAGES AND MISC CONFIGURATIONS
%-------------------------------------------------------
\usepackage{graphicx} 	     % Required for including images
\graphicspath{{figures/}{./}} % Figures path

\setlength{\parindent}{0pt} % Paragraph indentation
\setlength{\parskip}{10pt}  % Vertical space between paragraphs

\usepackage[bottom, hang]{footmisc} % Force footnotes to the bottom of the page and to the left margin
\setlength{\footnotemargin}{6pt}    % Horizontal space between the footnote marker and text

\usepackage[document]{ragged2e} % Left align all text in the document (i.e. have a ragged right margin)
\usepackage{lastpage} 			% Required to determine the total number of pages

\AtBeginEnvironment{quote}{\itshape\rmfamily} % Style quotations

\usepackage{appendix}  % Required for the appendices section
\usepackage{float} 		% Allows more precisely positioning floats e.g. \begin{figure}[H]
	

%-------------------------------------------------------
%	MARGINS
%-------------------------------------------------------
\usepackage[
	top=3cm, 	  			% Top margin
	bottom=3cm,  			% Bottom margin
	left=2cm,	  			% Left margin
	right=7.25cm, 			% Right margin
	marginparwidth=4.25cm,  % Width of the large right margin (sidebar) content
	marginparsep=1cm, 		% Padding between the main text and the large right margin (sidebar)
	footskip=1.5cm, 		% Space from the bottom margin to the baseline of the footer
	headsep=0.8cm, 			% Space from the top margin to the baseline of the header
	headheight=1cm, 		% Height of the header
	%showframe % Debugging frames, uncomment to show frames around margins
]{geometry}
\usepackage[oneside]{sidenotes} % Required for outputting content to the margin\

% Restyle \sidenote command from sidenotes package to reduce the font size
\makeatletter
\RenewDocumentCommand\sidenotetext{o o +m}{%      
	\IfNoValueOrEmptyTF{#1}{%
		\@sidenotes@placemarginal{#2}{\textsuperscript{\thesidenote}\footnotesize#3}%
		\refstepcounter{sidenote}%
	}{%
		\@sidenotes@placemarginal{#2}{\textsuperscript{#1}\footnotesize#3}%
	}%
}
\makeatother
% Environment to output to the full page width (excluding page margins)
\newenvironment{fullwidth}{ 
	% Pull into the right margin by the margin width and margin skip width
	\begin{adjustwidth}{}{\dimexpr-\marginparwidth-\marginparsep\relax} 
}{
	\end{adjustwidth}
}
% Environment to output content at two-thirds of the whole page width
\newenvironment{twothirdswidth}{ 
	% Pull into the right margin by half the margin width and half the margin skip width
	\begin{adjustwidth}{}{\dimexpr-0.5\marginparwidth-0.5\marginparsep\relax} 
}{
	\end{adjustwidth}
}

% Calculate the full text width (including the large right margin and margin padding)
\newlength{\fulltextwidth}
\setlength{\fulltextwidth}{\dimexpr\textwidth+\marginparwidth+\marginparsep\relax}
% New command to output an unnumbered sidenote, the optional argument is a vertical offset and the required parameter is the sidenote text
\newcommand{\nonumsidenote}[2][]{\sidenote[\hspace{0pt}][#1]{#2}} 
% New command to output a quotation sidenote, the optional argument is a vertical offset and the required parameter is the quotation text
\newcommand{\sidenotequote}[2][]{\sidenote[\hspace{0pt}][#1]{\rmfamily\itshape#2}} 
%-------------------------------------------------------
%	FONTS
%-------------------------------------------------------
\usepackage[utf8]{inputenc} 	  % Required for inputting international characters
\usepackage[T1]{fontenc} 		  % Output font encoding for international characters
\usepackage{mathpazo} 			  % Use Palatino for math
\usepackage[sfdefault]{plex-sans} % Use IBM Plex Sans as the sans font and make it the default
\usepackage{plex-serif} 		  % Use IBM Plex Serif as the serif font
\usepackage{plex-mono} 			  % Use IBM Plex Mono as the mono font

% ----------Rule Definations---------- {el}{l}{sb}
\newcommand{\textel}[1]{{\fontseries{el}\selectfont #1}} % ExtraLight weight command
\newcommand{\textl}[1]{{\fontseries{l}\selectfont #1}}   % Light weight
\newcommand{\textsb}[1]{{\fontseries{sb}\selectfont #1}} % SemiBold weight

%-------------------------------------------------------
%	HEADERS AND FOOTERS
%-------------------------------------------------------
\usepackage{fancyhdr} % Required for customizing headers and footers
\usepackage{multicol} % Used for multipule colums
\pagestyle{fancy} 	  % Enable custom headers and footers

\renewcommand{\headrulewidth}{0.5pt} % Top horizontal rule thickness

\fancyhf{} 			  % Clear default headers/footers

\fancyhead[L]{\reporttitle} 	   % Left header
\fancyhead[C]{} 				   % Center header
\fancyhead[R]{\rightheadercontent} % Right header

\fancyfoot[L]{} 							  % Left footer
\fancyfoot[C]{\thepage~of \pageref{LastPage}} % Center footer
\fancyfoot[R]{} 							  % Right footer

\usepackage{calc} 	  % Allows arithmetic in LaTeX commands
\newlength{\marginoffset}
\setlength{\marginoffset}{\marginparwidth + \marginparsep} % Calculate width of margin and buffer space before margin
\fancyheadoffset[R]{\marginoffset} 						   % Offset the header to account for the margin width
\fancyfootoffset[R]{\marginoffset} 						   % Offset the footer to account for the margin width

%-------------------------------------------------------
%	SECTIONS
%-------------------------------------------------------
\usepackage{titlesec} 		    % Required for modifying sections
\iftoggle{unnumberedsections}{  % Conditional logic for the unnumbered sections class options
	\setcounter{secnumdepth}{0} % Don't number sections at any level
}{
	\setcounter{secnumdepth}{3} % Number sections down to subsubsections
}
\titleformat
	{\section} % Section type being modified
	[block]    % Section layout type, can be: hang, block, display, runin, leftmargin, rightmargin, drop, wrap, frame
	{\raggedright\LARGE\bfseries} % Text formatting of the whole section, i.e. label and title
	{\thesection} % Section label (e.g. number) and its formatting
	{5pt} 		  % Horizontal space between the section label and title
	{} % Code before the section title
	[] % Code after the section title

\titlespacing*{\section}{0pt}{\baselineskip}{4pt} % Spacing around section titles, the order is: left, before and after

%------------------------------------------------

\titleformat
	{\subsection} % Section type being modified
	[block] 	  % Section layout type, can be: hang, block, display, runin, leftmargin, rightmargin, drop, wrap, frame
	{\raggedright\large\bfseries} % Text formatting of the whole section, i.e. label and title
	{\thesubsection} % Section label (e.g. number) and its formatting
	{5pt} 			 % Horizontal space between the section label and title
	{} % Code before the section title
	[] % Code after the section title

\titlespacing*{\subsection}{0pt}{0.75\baselineskip}{4pt} % Spacing around section titles, the order is: left, before and after

%------------------------------------------------

\titleformat
	{\subsubsection} % Section type being modified
	[block] 		 % Section layout type, can be: hang, block, display, runin, leftmargin, rightmargin, drop, wrap, frame
	{\raggedright\bfseries\normalsize} % Text formatting of the whole section, i.e. label and title
	{\thesubsubsection} % Section label (e.g. number) and its formatting
	{5pt} 				% Horizontal space between the section label and title
	{} % Code before the section title
	[] % Code after the section title

\titlespacing*{\subsubsection}{0pt}{0.5\baselineskip}{0pt} % Spacing around section titles, the order is: left, before and after

%------------------------------------------------

\titleformat
	{\paragraph} % Section type being modified
	[runin] 	 % Section layout type, can be: hang, block, display, runin, leftmargin, rightmargin, drop, wrap, frame
	{\bfseries}  % Text formatting of the whole section, i.e. label and title
	{} 			 % Section label (e.g. number) and its formatting
	{0pt} 		 % Horizontal space between the section label and title
	{} 			 % Code before the section title
	[~~~~] 		 % Code after the section title

\titlespacing*{\paragraph}{0pt}{0pt}{0pt} % Spacing around section titles, the order is: left, before and after

%-------------------------------------------------------
%	TABLE OF CONTENTS
%-------------------------------------------------------
\usepackage{titletoc}    % Required for manipulating the table of contents
\setcounter{tocdepth}{2} % Show entries in the table of contents down to subsections

\renewcommand{\contentsname}{Table of Contents} % Rename the table of contents section heading

\titlecontents{section} % Section type being modified
	[0pt] 				% Left indentation
	{\vspace{-2pt}} 	% Before code
	{\textsb{\thecontentslabel}~~\textsb} % Formatting of numbered sections of this type
	{\textsb} 			% Formatting of numberless sections of this type
	{\titlerule*[6pt]{.}\textsb{\contentspage}} % Formatting of the filler to the right of the heading and the page number
	[\vspace{2pt}] 		% After code

%------------------------------------------------

\titlecontents{subsection} % Section type being modified
	[0.5cm] 			   % Left indentation
	{\vspace{-8pt}} 	   % Before code
	{\textsb{\thecontentslabel}~~} % Formatting of numbered sections of this type
	{} 					   % Formatting of numberless sections of this type
	{\titlerule*[6pt]{.}\contentspage} % Formatting of the filler to the right of the heading and the page number
	[] 					   % After code
	
%------------------------------------------------

\titlecontents{subsubsection} % Section type being modified
	[1.35cm] 				  % Left indentation
	{\vspace{-4pt}} 		  % Before code
	{\textsb{\thecontentslabel}~~} % Formatting of numbered sections of this type
	{} 						  % Formatting of numberless sections of this type
	{\titlerule*[6pt]{.}\contentspage} % Formatting of the filler to the right of the heading and the page number
	[] 						  % After code

%-------------------------------------------------------
%	COLORS
%-------------------------------------------------------
\usepackage[usenames, svgnames]{xcolor,colortbl} % Required for defining and using custom colors
%\definecolor{cyan}{RGB}{0, 255, 255} % Define custom colors in this format

%-------------------------------------------------------
%	TABLES
%-------------------------------------------------------
\usepackage{booktabs}   % Required for better horizontal rules in tables (\toprule/bottomrule/addlinespace)
\usepackage{rotating}   % Rotating table labels
\usepackage{array} 	    % Required for manipulating table columns
\usepackage{tabularray} % LaTeX3, used for colspec / rowspec
\renewcommand{\arraystretch}{1.25} % Increase the height of table rows

% ----------Rule Definations----------
\newcolumntype{R}[1]{>{\raggedleft\arraybackslash}p{#1}}  % New right-aligned paragraph column type
\newcolumntype{L}[1]{>{\raggedright\arraybackslash}p{#1}} % New left-aligned (no justification) paragraph column type
\newcolumntype{C}[1]{>{\centering\arraybackslash}p{#1}}   % New centered paragraph column type

%-------------------------------------------------------
%	CAPTIONS
%-------------------------------------------------------
\usepackage{caption} 		   % Required for customizing captions
\captionsetup{skip=6pt} 	   % Vertical whitespace between figures/tables and the caption (default is 10pt)
\captionsetup{font={bf,small}} % Define caption font style
%\captionsetup{justification=raggedright, singlelinecheck=false} % Left-align captions

% The sidenotes package defines 4 custom caption styles for wide tables, wide figures, margin tables and margin figures
% To style these the same as in-text tables and figures, set these 4 styles to the same styling values
\DeclareCaptionStyle{widetable}{font={bf, footnotesize}, justification=raggedright, singlelinecheck=false}
\DeclareCaptionStyle{widefigure}{font={bf, footnotesize}, justification=raggedright, singlelinecheck=false}
\DeclareCaptionStyle{marginfigure}{font={bf, scriptsize}, justification=raggedright, singlelinecheck=false}
\DeclareCaptionStyle{margintable}{font={bf, scriptsize}, justification=raggedright, singlelinecheck=false}

%-------------------------------------------------------
%	LISTS
%-------------------------------------------------------
\usepackage{enumitem} 			  % Required for list customization

\setlist{itemsep=0pt, topsep=0pt} % Customize spacing around and inside lists

\renewcommand{\labelenumi}{\arabic{enumi}.}    % Customize the first level numbered list characters
\renewcommand{\labelenumii}{\alph{enumii}.}    % Customize the second level numbered list characters
\renewcommand{\labelenumiii}{\roman{enumiii}.} % Customize the third level numbered list characters

\renewcommand{\labelitemi}{\Large\raisebox{-1.5pt}{\textbullet}} % Customize the first level bullet point character
\renewcommand{\labelitemii}{\Large\raisebox{-1.5pt}{--}} 		 % Customize the second level bullet point character
\renewcommand{\labelitemiii}{\Large\raisebox{-1.5pt}{-}} 		 % Customize the third level bullet point character

%-------------------------------------------------------
%	LINKS
%-------------------------------------------------------
\usepackage{hyperref} % Required for links

\hypersetup{
	colorlinks=true, % Whether to color the text of links
	urlcolor=black,  % Color for \url and \href links
	linkcolor=black, % Color for \nameref links
	citecolor=black, % Color of reference citations
}

%-------------------------------------------------------
%	CODE LISTINGS
%-------------------------------------------------------
\usepackage{listings} % Required for formatted code blocks

\lstset{
	numbers=left, 		  		% Line numbers location
	numbersep=6pt,		  		% Whitespace to the right of line numbers
	numberstyle=\tiny\it, 		% Line numbers font style
	frame=single, 		  		% Border
	basicstyle=\tt\scriptsize,  % Code font style
	breaklines=true,      		% Break long lines into multiple lines automatically
	aboveskip=1.5\baselineskip, % Vertical whitespace before the listing
	tabsize=3, 			  		% How many space to a tab
}

%-------------------------------------------------------
%	CUSTOM COMMANDS
%-------------------------------------------------------
\newcommand{\reporttitle}[1]{\renewcommand{\reporttitle}{#1}}
\newcommand{\reportsubtitle}[1]{\renewcommand{\reportsubtitle}{#1}}
\newcommand{\reportauthors}[1]{\renewcommand{\reportauthors}{#1}}
\newcommand{\reportdate}[1]{\renewcommand{\reportdate}{#1}}
\newcommand{\rightheadercontent}[1]{\renewcommand{\rightheadercontent}{#1}}

%-------------------------------------------------------
%	BIBLIOGRAPHY
%-------------------------------------------------------
\usepackage[
	backend=biber,     % Use the biber backend for compiling the bibliography
	citestyle=numeric, % In-text citation style
	bibstyle=numeric,  % Bibliography style
]{biblatex}

\newcommand{\sidecite}[2][0pt]{\sidenote[][#1]{\textcite{#2}}}     % Parenthetical citations in the margin
\newcommand{\fullsidecite}[2][0pt]{\sidenote[][#1]{\fullcite{#2}}} % Full citations in the margin

\addbibresource{sample.bib} % BibLaTeX bibliography file

%------------------------------------------------
%	REPORT INFORMATION
%------------------------------------------------
\reporttitle{CompTIA Security+} 			    % The report title to appear on the title page and page headers
\reportsubtitle{Personal Notes {\color{red} Exam Objectives}} % Report subtitle, include new lines if needed
\reportauthors{\textit{Author:}\\\smallskip Scott \textit{Skrobel}} 		% Report authors/group/department, include new lines if needed
\reportdate{\today}														    % Report date, include new lines for additional information if needed
\rightheadercontent{\includegraphics[width=3cm]{logo.png}} 					% The content in the right header, leave empty for no right header content

%%%%%%%%%%%%%%%%%%%%%%%%%%%%%%%%%%%%%%%%%%%%%%%%%%%%%%%%%%%%%%%%
%							START DOC
%%%%%%%%%%%%%%%%%%%%%%%%%%%%%%%%%%%%%%%%%%%%%%%%%%%%%%%%%%%%%%%%
\begin{document}
%------------------------------------------------
%	Title Page
%------------------------------------------------
\thispagestyle{empty} % Suppress headers and footers on this page

\begin{fullwidth}     % Use the whole page width
	\vspace*{-0.075\textheight} 	           % Pull logo into the top margin
	%\hfill\includegraphics[width=5cm]{logo.png} % Company logo
	\vspace{0.15\textheight} 				   % Vertical whitespace
	% Report title, intentionally at less than full width for nice wrapping. Adjust width of \parbox and the font size as needed for your title to look good.
	\parbox{1.9\fulltextwidth}{\fontsize{50pt}{52pt}\selectfont\raggedright\textbf{\reporttitle}\par} 
	
	\vspace{0.03\textheight} 					  % Vertical whitespace	
	{\LARGE\textit{\textbf{\reportsubtitle}}\par} % Subtitle
	\vspace{0.04\textheight} 					  % Vertical whitespace
	{\color{black} \rule{\linewidth}{0.5mm}}      % Color Horizontal Line
	\vfill 										   % Vertical whitespace
	
	{\Large\reportauthors\par} % Report authors, group or department

	\vfill\vfill\vfill 			  % Vertical whitespace
	\begin{flushright}
		
	{\large\reportdate\par}    % Report date
	\end{flushright}
\end{fullwidth}

\newpage
%------------------------------------------------
%	TABLE OF CONTENTS
%------------------------------------------------

\begin{twothirdswidth} % Content in this environment to be at two-thirds of the whole page width
	\tableofcontents   % Output the table of contents, automatically generated from the section commands used in the document
\end{twothirdswidth}

\newpage

%------------------------------------------------
%	SECTIONS
%------------------------------------------------
\section{Threats, Attacks and Vulnerabilities}



\subsection{Compare and contrast types of attacks}
\begin{fullwidth}
    \begin{description}\itemsep2pt
        \item[Phishing]: Social engineering tactic to acquire personal information from a fake email with a clickable link.
        \item[Smishing]: (SMS phishing) The use of deceptive text messages into divulging sensitive information.
        \item[Vishing]: (Voice Phishing) Impersonates a trusted entity, such as a bank to trick into giving information.
        \item[Spam]: Unsolicited inappropriate messages sent with the purpose of spreading malware, advertising, phishing. 
        \item[Spear phishing]: Targeted type of phishing attack to make the scam convincing, often with insider information.
        \item[Dumpster diving]: Searching through an organization's or individual's trash to find sensitive information.
        \item[Shoulder surfing]: Observing a victim’s screen or keyboard to obtain sensitive information.
        \item[Pharming]: Manipulating the DNS system to resolve fake domain names, to lead them to a fake website.
        \item[Tailgating]: Physical security breach by following authorized person to get access to secure areas.
        \item[Eliciting information]: Psychological tactics to encourage individuals to share their knowledge willingly. 
        \item[Whaling]: Spear phishing for high-profile executives in an organization.
        \item[Prepending]: Organizing, manipulating, and structuring data in various applications.
        \item[Identity fraud]: An individual wrongfully obtains and uses someone else's personal data in a deceptive manner.
        \item[Invoice scams]: Impersonation of a legit business to deceive individuals into paying fraudulent invoices. 
        \item[Credential harvesting]: Tricking someone into disclosing login credentials to access sensitive info.
        \item[Reconnaissance]: Initial phase to gather intelligence via passive and active techniques.
        \item[Hoax]: Fabrication intended to deceive or trick individuals into believing false information or events.
        \item[Impersonation]: Masquerading as a legitimate user or entity to gain unauthorized access to information.
        \item[Watering hole attack]: Infecting a commonly visited website of a targeted specific group.
        \item[Typosquatting]: URL hijacking or domain squatting of fake domains which resemble a legit website.
        \item[Pretexting]: Fabricated scenario involving direct interaction to obtain sensitive information.
        \item[Influence campaigns]: Coordinated effort to shape public opinion, influence perceptions, and manipulate.
        \begin{itemize}\itemsep2pt
            \item[]\textbf{Hybrid warfare}: Blends conventional warfare tactics with unconventional method
        \end{itemize}
        \item[Principles] (reasons for effectiveness)
        \begin{itemize}\itemsep2pt
            \item \textbf{Authority}: The actor acts as an individual of authority
            \item \textbf{Intimidation}: Frightening or threatening the victim.
            \item \textbf{Consensus}: Convince based on what’s normally expected.
            \item \textbf{Scarcity}: Limited resources and time to act.
            \item \textbf{Familiarity}: The victim is well known.
            \item \textbf{Trust}: Gain their confidence, be their friend.
            \item \textbf{Urgency}: Limited time to act, rush the victim.
        \end{itemize}
    \end{description}
\end{fullwidth}

\subsection{Given a scenario, analyze indicators of compromise}
\subsubsection*{\color{green}Malware}
\begin{fullwidth}
    \begin{description}\itemsep2pt
        \item[Ransomware]: Denies access to a computer system or data until a ransom is paid.
        \item[Trojan]: A form of malware that pretends to be a harmless application.
        \item[Worm]: A self-contained infection that can spread itself through networks, emails, and messages.
        \item[PUP's]: Potentially Unwanted Programs software applications that may exhibit undesirable characteristics.
        \item[Memory-resident malware]: Operates primarily in a computer's volatile memory (RAM) rather than with files
        \item[Command and control]: (C2) Centralized server used by attackers to manager compromised devices.
        \item[Bots]: AI inside an infected machine performs specific actions as a part of a larger entity known as a botnet.
        \item[Cryptomalware]: A malicious program that encrypts programs and files on the computer to extort money.
        \item[Logic Bomb]: A malicious program that lies dormant until a specific date or event occurs.
        \item[Spyware]: Software that installs itself to spy and sends stolen info back to the host machine.
        \item[Keyloggers]: A malicious program that saves all of the keystrokes of the infected machine.
        \item[Remote Access Trojan]: (RAT) A remotely operated Trojan.
        \item[Rootkit]: A backdoor program that allows full remote access to a system.
        \item[Backdoor]: Allows for full access to a system remotely.
    \end{description}
\end{fullwidth}

\subsubsection*{\color{blue}Password Attacks}
\begin{fullwidth}
    \begin{description}
        \item[Sprying] a
    \end{description}
\end{fullwidth}

\subsubsection*{\color{red}Physical attacks}
\begin{fullwidth}
    \begin{description}
        \item[Sprying] a
    \end{description}
\end{fullwidth}

\section{Architecture and Design}








\end{document}